\chapter{Lesson 2: Shopping}

This lesson covers buying things: asking prices, reading large numbers in won, asking a
shopkeeper for things, and saying what you want with \ko{\til 고 싶다}. In class we
role-played a market stall with paper money, which is where most of the example
sentences below come from.

\section{Vocabulary}

\begin{table}[H]
\centering
\begin{tabular}{lll}
\toprule
Korean & RR & English \\
\midrule
\ko{가게} & gage & shop, store \\
\ko{시장} & sijang & market \\
\ko{백화점} & baekhwajeom & department store \\
\ko{편의점} & pyeonuijeom & convenience store \\
\ko{옷} & ot & clothes \\
\ko{신발} & sinbal & shoes \\
\ko{가방} & gabang & bag \\
\ko{모자} & moja & hat, cap \\
\ko{선물} & seonmul & gift \\
\ko{값}/\ko{가격} & gap / gagyeok & price \\
\ko{사다} & sada & to buy \\
\ko{팔다} & palda & to sell \\
\ko{싸다} & ssada & to be cheap \\
\ko{비싸다} & bissada & to be expensive \\
\ko{크다} / \ko{작다} & keuda / jakda & to be big / small \\
\ko{현금} & hyeongeum & cash \\
\ko{카드} & kadeu & (credit) card \\
\ko{영수증} & yeongsujeung & receipt \\
\ko{색깔} & saekkkal & colour \\
\bottomrule
\end{tabular}
\end{table}

\subsection{Useful phrases in a shop}

\begin{table}[H]
\centering
\begin{tabular}{lll}
\toprule
Korean & RR & English \\
\midrule
\ko{어서 오세요.} & Eoseo oseyo. & Welcome. (said by the shopkeeper) \\
\ko{뭐 찾으세요?} & Mwo chajeuseyo? & What are you looking for? \\
\ko{이거 얼마예요?} & Igeo eolmayeyo? & How much is this? \\
\ko{좀 깎아 주세요.} & Jom kkakka juseyo. & Please give me a discount. \\
\ko{다른 색깔 있어요?} & Dareun saekkkal isseoyo? & Do you have another colour? \\
\ko{이걸로 주세요.} & Igeollo juseyo. & I'll take this one. \\
\ko{카드 돼요?} & Kadeu dwaeyo? & Can I pay by card? \\
\ko{여기 있어요.} & Yeogi isseoyo. & Here you are. \\
\ko{안녕히 가세요.} & Annyeonghi gaseyo. & Goodbye. (to someone leaving) \\
\bottomrule
\end{tabular}
\end{table}

\section{Demonstratives: \ko{이}, \ko{그}, \ko{저}}

Korean distinguishes three distances. In speech, \ko{것} (thing) is usually shortened to
\ko{거}.

\begin{center}
\begin{tikzpicture}[every node/.style={font=\small}]
  \node[draw, circle, fill=blue!10, minimum size=0.9cm] (sp) at (0,0) {S};
  \node[draw, circle, fill=green!10, minimum size=0.9cm] (li) at (4,0) {L};
  \node[font=\scriptsize] at (0,-0.75) {speaker};
  \node[font=\scriptsize] at (4,-0.75) {listener};
  \node[draw, rounded corners, fill=orange!15] (a) at (-0.7,1.4) {\ko{이거} igeo};
  \node[draw, rounded corners, fill=orange!15] (b) at (4.9,1.3) {\ko{그거} geugeo};
  \node[draw, rounded corners, fill=orange!15] (c) at (2,3.0) {\ko{저거} jeogeo};
  \draw[dashed] (sp) -- (a); \draw[dashed] (li) -- (b);
  \draw[dashed] (sp) -- (c); \draw[dashed] (li) -- (c);
  \node[font=\scriptsize, align=left, anchor=west] at (6.2,1.3) {\ko{이} near the speaker\\ \ko{그} near the listener\\ \quad (or already mentioned)\\ \ko{저} far from both};
\end{tikzpicture}
\end{center}

\begin{table}[H]
\centering
\begin{tabular}{llll}
\toprule
 & \ko{이} (this) & \ko{그} (that) & \ko{저} (that over there) \\
\midrule
+ noun & \ko{이 가방} & \ko{그 가방} & \ko{저 가방} \\
thing & \ko{이것 / 이거} & \ko{그것 / 그거} & \ko{저것 / 저거} \\
thing + \ko{이} (subject) & \ko{이게} & \ko{그게} & \ko{저게} \\
thing + \ko{을} (object) & \ko{이걸} & \ko{그걸} & \ko{저걸} \\
place & \ko{여기} (here) & \ko{거기} (there) & \ko{저기} (over there) \\
\bottomrule
\end{tabular}
\end{table}

\section{Grammar 1: Prices and large numbers}

Money always uses Sino-Korean numbers with \ko{원} won. The key difference from English
is that large numbers are grouped by \textbf{four digits} (units of \ko{만}, 10\,000),
not three.

\begin{center}
\begin{tikzpicture}[every node/.style={font=\small}]
  \foreach \i/\d/\u in {0/1/\ko{천만}, 1/2/\ko{백만}, 2/3/\ko{십만}, 3/4/\ko{만}, 4/5/\ko{천}, 5/6/\ko{백}, 6/7/\ko{십}, 7/8/\ko{일}} {
    \node[draw, minimum width=1cm, minimum height=0.8cm] at (\i*1.1,0) {\d};
    \node[font=\scriptsize] at (\i*1.1,0.65) {\u};
  }
  \draw[thick, blue!60!black] (-0.5,-0.55) -- (3.8,-0.55);
  \draw[thick, blue!60!black] (-0.5,-0.45) -- (-0.5,-0.65); \draw[thick, blue!60!black] (3.8,-0.45) -- (3.8,-0.65);
  \draw[thick, red!60!black] (3.9,-0.55) -- (8.2,-0.55);
  \draw[thick, red!60!black] (3.9,-0.45) -- (3.9,-0.65); \draw[thick, red!60!black] (8.2,-0.45) -- (8.2,-0.65);
  \node[font=\scriptsize, blue!60!black] at (1.65,-0.9) {\ko{천이백삼십사만} (1234 \ko{만})};
  \node[font=\scriptsize, red!60!black] at (6.05,-0.9) {\ko{오천육백칠십팔}};
  \node[font=\scriptsize] at (4,-1.45) {12,345,678 $=$ \ko{천이백삼십사만 오천육백칠십팔}};
\end{tikzpicture}
\end{center}

\begin{important}[Reading prices]
\begin{itemize}
  \item Do not say \ko{일} before \ko{십}, \ko{백}, \ko{천} or \ko{만}: 10\,000 is
        \ko{만} man, not \ko{일만}; 1\,500 is \ko{천오백} cheonobaek.
  \item Split the number into groups of four from the right, read each group, and attach
        \ko{만} to the upper group: 45\,000 = 4\,{\textbar}\,5000 $\to$ \ko{사만 오천 원}
        saman ocheon won.
  \item Common prices: 1\,000 \ko{천 원} cheon won, 5\,000 \ko{오천 원} ocheon won,
        10\,000 \ko{만 원} man won, 50\,000 \ko{오만 원} oman won (the largest banknote).
\end{itemize}
\end{important}

\begin{example}[Prices]
Read aloud in the role-play:
\begin{itemize}
  \item \ko{3,500원} \ko{삼천오백 원} samcheonobaek won (a coffee)
  \item \ko{12,000원} \ko{만 이천 원} man icheon won (a lunch set)
  \item \ko{89,000원} \ko{팔만 구천 원} palman gucheon won (a pair of shoes)
  \item \ko{150,000원} \ko{십오만 원} siboman won (a winter coat): ``fifteen \ko{만}'',
        pronounced [\ko{시보만}] by liaison.
\end{itemize}
\end{example}

\subsection{Quantity: noun + number + counter}

When buying several things, the order is \emph{noun} (+ particle) + \emph{native number} +
\emph{counter}.

\begin{example}[Ordering quantities]
At the fruit stand and the caf\'e:
\begin{itemize}
  \item \ko{사과 다섯 개 주세요.} Sagwa daseot gae juseyo. (Five apples, please.)
  \item \ko{콜라 두 병하고 물 한 병 주세요.} Kolla du byeonghago mul han byeong juseyo.
        (Two bottles of cola and one bottle of water, please.)
  \item \ko{아메리카노 한 잔에 얼마예요?} Amerikano han jane eolmayeyo? (How much is one
        Americano?) --- \ko{에} here means ``per''.
  \item \ko{이 공책은 한 권에 천 원이에요.} I gongchaegeun han gwone cheon wonieyo. (This
        notebook is 1,000 won each.)
\end{itemize}
\end{example}

\section{Grammar 2: \ko{\til 아/어 주세요} (please do something for me)}

\ko{주세요} alone means ``please give (me)'': \ko{물 주세요} (water, please). Attached to
the \ko{\til 아/어} form of a verb (the \ko{해요} form without \ko{요}), it asks someone to
do that action for you.

\begin{table}[H]
\centering
\begin{tabular}{llll}
\toprule
Verb & \ko{\til 아/어 주세요} & RR & Meaning \\
\midrule
\ko{깎다} (cut down) & \ko{깎아 주세요} & kkakka juseyo & give me a discount \\
\ko{보여 주다} (show) & \ko{보여 주세요} & boyeo juseyo & please show me \\
\ko{바꾸다} (exchange) & \ko{바꿔 주세요} & bakkwo juseyo & please exchange it \\
\ko{싸다} (wrap) & \ko{싸 주세요} & ssa juseyo & please wrap it \\
\ko{포장하다} (pack) & \ko{포장해 주세요} & pojanghae juseyo & to go, please \\
\bottomrule
\end{tabular}
\end{table}

\begin{remark}
\ko{싸다} is a homonym: ``to be cheap'' and ``to wrap''. Context sorts it out:
\ko{싸요!} (It's cheap!) versus \ko{싸 주세요} (Wrap it, please).
\end{remark}

\section{Grammar 3: \ko{\til 고 싶다} (want to)}

To say you want to \emph{do} something, attach \ko{\til 고 싶다} to the verb stem (no
vowel harmony; the stem is unchanged). \ko{싶다} then conjugates as a normal adjective:
\ko{싶어요}, \ko{싶었어요}.

\begin{example}[Wants]
Stem + \ko{고 싶어요}:
\begin{itemize}
  \item \ko{새 신발을 사고 싶어요.} Sae sinbareul sago sipeoyo. (I want to buy new
        shoes.)
  \item \ko{뭐 먹고 싶어요? --- 떡볶이를 먹고 싶어요.} Mwo meokgo sipeoyo? ---
        Tteokbokkireul meokgo sipeoyo. (What do you want to eat? --- I want to eat
        tteokbokki.)
  \item \ko{어제 영화를 보고 싶었어요. 그런데 시간이 없었어요.} Eoje yeonghwareul bogo
        sipeosseoyo. Geureonde sigani eopseosseoyo. (Yesterday I wanted to see a movie,
        but I didn't have time.)
  \item \ko{아무것도 하고 싶지 않아요.} Amugeotdo hago sipji anayo. (I don't want to do
        anything.)
\end{itemize}
\end{example}

\begin{important}[First vs.\ third person]
\ko{\til 고 싶다} describes the \emph{speaker's} (or, in a question, the listener's)
feelings, which only they can know directly. For a third person, use
\ko{\til 고 싶어하다}: \ko{동생이 강아지를 키우고 싶어해요.} Dongsaengi gangajireul kiugo
sipeohaeyo. (My little brother wants to raise a puppy.) The same rule applies to other
feeling adjectives: \ko{좋다} $\to$ \ko{좋아하다}.
\end{important}

To want a \emph{thing} (a noun), Korean prefers ``noun + \ko{을/를} + verb + \ko{고 싶다}''
(\ko{커피를 마시고 싶어요}) or noun + \ko{이/가 필요해요} (I need ...), rather than a direct
``I want X''.

\section{Dialogue practice (paraphrased from class role-play)}

\begin{quote}
\ko{점원: 어서 오세요. 뭐 찾으세요?}\\
\ko{손님: 가방을 사고 싶어요. 저 검은색 가방 좀 보여 주세요.}\\
\ko{점원: 여기 있어요. 요즘 인기가 많아요.}\\
\ko{손님: 얼마예요?}\\
\ko{점원: 육만 오천 원이에요.}\\
\ko{손님: 좀 비싸요. 좀 깎아 주세요.}\\
\ko{점원: 그럼 육만 원에 드릴게요.}\\
\ko{손님: 좋아요. 카드 돼요?}
\end{quote}

\noindent Clerk: Welcome, what are you looking for? / Customer: I want to buy a bag.
Please show me that black bag. / Clerk: Here it is. It's popular these days. / Customer:
How much is it? / Clerk: 65,000 won. / Customer: That's a bit expensive. Please give me a
discount. / Clerk: Then I'll give it to you for 60,000 won. / Customer: Great. Can I pay
by card?

\begin{aside}
\textbf{Markets and bargaining.} Traditional markets such as Namdaemun and Dongdaemun in
Seoul are places where gentle bargaining is expected; department stores and convenience
stores have fixed prices. Vendors at markets often throw in a little extra for free,
called \ko{덤} deom, as a sign of \ko{정} jeong (warm-hearted affection, a very important
concept in Korean culture). There is no tipping in Korea, and the price on the tag
includes tax.
\end{aside}

\section*{Review summary}
\begin{itemize}
  \item Prices: Sino-Korean numbers + \ko{원}; group digits by four (\ko{만}); no
        \ko{일} before \ko{십/백/천/만}.
  \item Quantities: noun + native number + counter (\ko{사과 세 개}, \ko{물 한 병}).
  \item \ko{이/그/저} = near me / near you / far from both; \ko{이거, 그거, 저거}.
  \item \ko{noun 주세요} (give me); \ko{verb-아/어 주세요} (do it for me).
  \item \ko{\til 고 싶다} = want to (1st/2nd person); \ko{\til 고 싶어하다} for others.
\end{itemize}
